\documentclass[11pt,a4paper]{article}

\usepackage[utf8]{inputenc}
\usepackage[T1]{fontenc}
\usepackage{lmodern}
\usepackage[spanish,es-noshorthands]{babel}
\usepackage{xcolor}
\usepackage[most]{tcolorbox}
\usepackage{listings}
\usepackage{fancyhdr}
\usepackage{lastpage}
\usepackage{enumitem}
\usepackage[hidelinks]{hyperref}

\setlength{\headheight}{30pt}
\pagestyle{fancy}
\fancyhf{}
\renewcommand{\headrulewidth}{0pt}
\fancyhead[R]{%
  \fontsize{8}{9.6}\selectfont
  Licenciatura en Ciencias de la Computación\\
  Sistemas Operativos\\[2pt]
  \rule{4.2cm}{0.5pt}
}
\fancyfoot[C]{\thepage\ / \pageref{LastPage}}

\newtcolorbox{infobox}{
  colback=blue!6!white, colframe=blue!45!white,
  arc=2mm, boxrule=0.6pt, left=8pt, right=8pt, top=6pt, bottom=6pt,
  title=INFO, fonttitle=\bfseries, coltitle=blue!45!black,
  colbacktitle=blue!6!white, boxsep=1pt, breakable
}

\newtcolorbox{extrabox}{
  colback=orange!8!white, colframe=orange!55!white,
  arc=2mm, boxrule=0.6pt, left=8pt, right=8pt, top=6pt, bottom=6pt,
  title=EXTRA, fonttitle=\bfseries, coltitle=orange!55!black,
  colbacktitle=orange!8!white, boxsep=1pt, breakable
}

\newtcolorbox{notebox}{
  colback=black!5!white, colframe=black!25!white,
  arc=2mm, boxrule=0.5pt, left=8pt, right=8pt, top=6pt, bottom=6pt,
  boxsep=1pt, breakable
}

% Caja nueva para los desafíos BONUS - POR PUNTOS de los tracks
\newtcolorbox{bonusbox}{
  colback=purple!6!white, colframe=purple!55!white,
  arc=2mm, boxrule=0.6pt, left=8pt, right=8pt, top=6pt, bottom=6pt,
  title=BONUS -- POR PUNTOS, fonttitle=\bfseries, coltitle=purple!55!black,
  colbacktitle=purple!6!white, boxsep=1pt, breakable
}

% "Medalla" con el puntaje de cada track, en la misma línea del título
\newtcbox{\trackbadge}[1][]{
  on line, arc=1.5pt, colframe=blue!35!black, colback=blue!35!black,
  boxrule=0pt, boxsep=1pt, left=6pt, right=6pt, top=3pt, bottom=3pt,
  fontupper=\bfseries\footnotesize\color{white}
}

\lstdefinestyle{codestyle}{
  backgroundcolor=\color{black!5},
  rulecolor=\color{black!25},
  frame=single,
  framerule=0.5pt,
  basicstyle=\ttfamily\small,
  breaklines=true,
  showstringspaces=false,
  tabsize=4,
  keywordstyle=\color{blue!60!black}\bfseries,
  commentstyle=\color{green!40!black},
  stringstyle=\color{orange!70!black},
  columns=fullflexible,
  upquote=true
}
\lstset{style=codestyle}
\setlist{nosep}

\title{Trabajo Práctico 4 -- Ejercicios extra}
\author{}
\date{}

\begin{document}
\maketitle
\thispagestyle{fancy}
\begin{center}
SISTEMAS OPERATIVOS
\end{center}
\newpage

\renewcommand{\contentsname}{Tracks}
\tableofcontents
\newpage

Estos tracks amplían el práctico de planificación con experimentos comparables.
Registrá versión del núcleo, cantidad de núcleos lógicos y carga del sistema
durante cada observación: los valores de CPU y de latencia no son portables
entre máquinas.

Los tracks son independientes entre sí: elegí uno después de completar el
práctico regular y no necesitás resolver otro track para empezar. Las cajas
EXTRA proponen una extensión opcional del track; las cajas BONUS -- POR PUNTOS
plantean una implementación más profunda. El práctico regular no requiere un
informe; para obtener puntos de bonus, entregá una explicación breve con
capturas y quedate disponible para una demostración.

Varios experimentos de este documento usan políticas de tiempo real o
afinidad de CPU. No los ejecutes sobre una máquina compartida ni contra
procesos del sistema, y asegurate de poder terminar cualquier tarea que crees
antes de dejar la terminal.

\section{Políticas de tiempo real: \texttt{SCHED\_FIFO} y \texttt{SCHED\_RR} con \texttt{chrt} (7/10)}
\trackbadge{TRACK - 7/10}

\emph{Descripción:} comparar el tiempo compartido normal con dos políticas de
tiempo real usando un mismo programa de carga.

Escribí un programa \texttt{carga.c} que ejecute un bucle de cómputo (por
ejemplo, incrementar un contador durante varios segundos) sin bloquearse por
E/S. Lanzá varias instancias bajo \texttt{SCHED\_OTHER}, \texttt{SCHED\_FIFO}
y \texttt{SCHED\_RR} usando \texttt{chrt}, y compará cuánta CPU recibe cada
una.

\begin{lstlisting}[language=bash]
$ gcc -Wall -Wextra -O0 -g carga.c -o carga
$ ./carga & pid1=$!
$ ./carga & pid2=$!
$ chrt -p $pid1
$ sudo chrt -f -p 10 $pid1
$ sudo chrt -r -p 10 $pid2
$ ps -o pid,cls,rtprio,ni,pri,pcpu,comm -p "$pid1,$pid2"
$ top -H
\end{lstlisting}

Repetí la observación con dos procesos \texttt{SCHED\_FIFO} de la misma
prioridad y con dos procesos \texttt{SCHED\_RR} de la misma prioridad.
Explicá por qué \texttt{SCHED\_FIFO} puede dejar sin CPU a un proceso de
tiempo compartido mientras el de tiempo real no se bloquea, y qué diferencia
notás frente a \texttt{SCHED\_RR} cuando hay más de un proceso de tiempo real
listo.

Terminá siempre los procesos de tiempo real con \texttt{kill} antes de cerrar
la terminal: un \texttt{SCHED\_FIFO} que nunca cede la CPU puede volver el
sistema poco responsivo.

\begin{extrabox}
Consultá \texttt{man 2 sched\_setscheduler} y explicá qué devuelve
\texttt{sched\_get\_priority\_min()} y \texttt{sched\_get\_priority\_max()}
para cada política. Relacioná el resultado con los valores que aceptó
\texttt{chrt} en tu experimento.
\end{extrabox}

\newpage
\section{Afinidad de CPU con \texttt{taskset} y \texttt{sched\_getaffinity()} (5/10)}
\trackbadge{TRACK - 5/10}

\emph{Descripción:} fijar un proceso a un subconjunto de CPUs lógicas y
observar el efecto sobre \texttt{PSR} y el uso por núcleo.

\begin{enumerate}
  \item Consultá cuántas CPUs lógicas tiene el sistema y la máscara de
  afinidad por defecto de un proceso.
  \item Lanzá dos procesos de carga (podés reusar \texttt{carga.c} del track
  anterior) sin restricción de afinidad y observá en qué CPU corre cada uno
  con la columna \texttt{PSR}.
  \item Fijá ambos procesos a la misma CPU con \texttt{taskset} y volvé a
  observar \texttt{\%CPU} y \texttt{PSR}.
  \item Fijá cada proceso a una CPU distinta y compará.
\end{enumerate}

\begin{lstlisting}[language=bash]
$ nproc
$ taskset -p $pid1
$ ps -o pid,psr,pcpu,comm -p "$pid1,$pid2"
$ taskset -cp 0 $pid1
$ taskset -cp 0 $pid2
$ ps -o pid,psr,pcpu,comm -p "$pid1,$pid2"
$ taskset -cp 1 $pid2
$ ps -o pid,psr,pcpu,comm -p "$pid1,$pid2"
\end{lstlisting}

Relacioná lo observado con el hecho de que dos procesos en la misma CPU deben
repartirse el tiempo del planificador, mientras que en CPUs distintas pueden
ejecutarse en paralelo. Explicá qué es una máscara de afinidad y por qué
reducirla no aumenta la prioridad de un proceso.

\begin{extrabox}
Desde el propio programa, usá \texttt{sched\_getaffinity()} y
\texttt{sched\_setaffinity()} para leer e imprimir la máscara de CPUs
disponibles al inicio. Compará el resultado con el de \texttt{taskset -p}
ejecutado externamente sobre el mismo PID.
\end{extrabox}

\newpage
\section{Límites de CPU con \texttt{cgroups} v2 (6/10)}
\trackbadge{TRACK - 6/10}

\emph{Descripción:} crear un grupo de control (\emph{cgroup}) y limitar la
cuota de CPU de los procesos que contiene.

Trabajá en una máquina de prueba donde tengas permisos sobre
\texttt{/sys/fs/cgroup}. Necesitás un sistema con \emph{cgroups} v2 montado;
comprobalo antes de continuar.

\begin{lstlisting}[language=bash]
$ cat /sys/fs/cgroup/cgroup.controllers
$ sudo mkdir /sys/fs/cgroup/practico-so
$ cat /sys/fs/cgroup/practico-so/cpu.max
$ ./carga & pid=$!
$ echo $pid | sudo tee /sys/fs/cgroup/practico-so/cgroup.procs
$ cat /sys/fs/cgroup/practico-so/cgroup.procs
$ top -p $pid
\end{lstlisting}

Con el proceso ya dentro del \emph{cgroup} y consumiendo CPU, limitá su cuota
a un 20\% aproximado de una CPU y observá cómo cambia \texttt{\%CPU} en
\texttt{top}. \texttt{cpu.max} acepta dos valores: la cuota en microsegundos
y el período; por ejemplo, \texttt{20000 100000} permite usar 20\,ms de cada
100\,ms.

\begin{lstlisting}[language=bash]
$ echo "20000 100000" | sudo tee /sys/fs/cgroup/practico-so/cpu.max
$ top -p $pid
$ echo "max 100000" | sudo tee /sys/fs/cgroup/practico-so/cpu.max
$ top -p $pid
\end{lstlisting}

Al terminar, sacá el proceso del \emph{cgroup}, matalo si sigue vivo y borrá
el directorio. No modifiques ni elimines \emph{cgroups} que no hayas creado
vos.

\begin{lstlisting}[language=bash]
$ kill $pid
$ sudo rmdir /sys/fs/cgroup/practico-so
\end{lstlisting}

Explicá la diferencia entre limitar la cuota de CPU de un \emph{cgroup} y
bajar la prioridad (\texttt{nice}) de un proceso: ¿en qué se parecen los
efectos observables y en qué se diferencia el mecanismo?

\begin{extrabox}
Poné dos procesos de carga en el mismo \emph{cgroup} y explorá
\texttt{cpu.weight} en lugar de \texttt{cpu.max}. Compará cómo reparte la CPU
entre ambos frente a dejarlos sin \emph{cgroup}.
\end{extrabox}

\newpage
\section{Trazar al planificador con \texttt{perf sched} (8/10)}
\trackbadge{TRACK - 8/10}

\emph{Descripción:} usar \texttt{perf} para registrar eventos de
planificación y relacionarlos con los contadores de \texttt{/proc} que ya
usaste en el práctico regular.

Necesitás permisos para usar \texttt{perf record}; en muchos sistemas de
prueba alcanza con \texttt{sudo} o con ajustar
\texttt{kernel.perf\_event\_paranoid} temporalmente. No modifiques ese
parámetro en una máquina compartida sin autorización.

\begin{lstlisting}[language=bash]
$ perf --version
$ sudo perf sched record -- ./carga
$ sudo perf sched latency
$ sudo perf sched map | head -50
\end{lstlisting}

\texttt{perf sched latency} resume, por tarea, cuánto tiempo esperó para
ejecutarse y cuántas veces fue planificada. \texttt{perf sched map} muestra,
en una línea de tiempo por CPU, cuándo cada tarea entra y sale del
procesador.

Repetí la traza mientras corren dos o más instancias de \texttt{carga} con
distinta niceness, y buscá en la salida los momentos en que una tarea
desaloja a otra. Relacioná lo que ves con \texttt{voluntary\_ctxt\_switches}
y \texttt{nonvoluntary\_ctxt\_switches} de \texttt{/proc/<pid>/status} para
las mismas tareas durante el mismo intervalo.

\begin{bonusbox}
Escribí un script que combine tres fuentes para un mismo experimento: el
resumen de \texttt{perf sched latency}, los contadores de
\texttt{/proc/<pid>/status} antes y después, y una tabla de \texttt{\%CPU}
tomada de \texttt{ps} a intervalos regulares.

El informe debe identificar al menos un caso de desalojo visible en
\texttt{perf sched map} y explicar, con los tres tipos de dato, por qué
ocurrió en ese momento y no antes.

Compará el resultado con el mismo experimento corrido con
\texttt{SCHED\_OTHER} y con \texttt{SCHED\_RR} y discutí qué cambia en la
frecuencia y el motivo de los desalojos.
\end{bonusbox}

\newpage
\section{\texttt{sched\_yield()} y costo de un cambio de contexto (4/10)}
\trackbadge{TRACK - 4/10}

\emph{Descripción:} medir cuánto tarda un cambio de contexto provocado
voluntariamente y compararlo con uno involuntario.

Escribí un programa \texttt{yield.c} que en un bucle llame a
\texttt{sched\_yield()} muchas veces (por ejemplo, un millón) y mida el
tiempo total con \texttt{clock\_gettime(CLOCK\_MONOTONIC, ...)}. Dividí el
tiempo total por la cantidad de llamadas para estimar un costo promedio.

\begin{lstlisting}[language=bash]
$ gcc -Wall -Wextra -O0 -g yield.c -o yield
$ ./yield
$ grep -E 'voluntary_ctxt_switches|nonvoluntary_ctxt_switches' /proc/self/status
\end{lstlisting}

Corré el mismo programa solo en el sistema y luego con otro proceso de carga
compitiendo por la misma CPU (usá \texttt{taskset} para forzarlos a la misma
CPU lógica). Compará el costo promedio de \texttt{sched\_yield()} en ambos
casos y contá cuántos cambios de contexto voluntarios e involuntarios
acumuló el proceso en \texttt{/proc/<pid>/status}.

Explicá por qué \texttt{sched\_yield()} no garantiza que otra tarea
específica reciba la CPU, y por qué el costo medido depende de si hay otra
tarea lista para ejecutarse en esa misma CPU.

\begin{extrabox}
Repetí la medición reservando el proceso a una sola CPU con \texttt{taskset
-c} y comparala contra dejarlo sin restricción de afinidad en una máquina con
varios núcleos libres.
\end{extrabox}

\newpage
\section{Prioridad dinámica, \texttt{nice} y \emph{starvation} (6/10)}
\trackbadge{TRACK - 6/10}

\emph{Descripción:} provocar un escenario de inanición (\emph{starvation})
leve entre procesos de distinta niceness y decidir si corresponde llamarlo
así.

Lanzá cuatro o más procesos de carga: la mayoría con \texttt{nice -n 19}
(mínima preferencia) y uno solo con \texttt{nice -n -5} o \texttt{nice -n 0},
todos compitiendo por la misma cantidad reducida de CPUs lógicas (podés usar
\texttt{taskset} para limitarlos a una o dos CPUs).

\begin{lstlisting}[language=bash]
$ taskset -c 0 nice -n 19 ./carga & p1=$!
$ taskset -c 0 nice -n 19 ./carga & p2=$!
$ taskset -c 0 nice -n 19 ./carga & p3=$!
$ taskset -c 0 nice -n -5 ./carga & p4=$!
$ ps -o pid,ni,pri,pcpu,comm -p "$p1,$p2,$p3,$p4"
$ watch -n 1 'ps -o pid,ni,pri,pcpu,comm -p "'"$p1,$p2,$p3,$p4"'"'
\end{lstlisting}

Registrá el \texttt{\%CPU} de cada proceso a lo largo de al menos un minuto.
Contestá: ¿alguno de los procesos con \texttt{nice 19} llegó a \texttt{\%CPU}
cercano a cero de forma sostenida, o todos recibieron algo de tiempo aunque
fuera poco? Investigá por qué \texttt{SCHED\_OTHER} en Linux está diseñado
para evitar la inanición total incluso con niceness muy dispar, a diferencia
de lo que podría ocurrir con una prioridad estática mal configurada en una
política de tiempo real.

\begin{extrabox}
Repetí el experimento reemplazando al proceso de \texttt{nice -5} por uno
\texttt{SCHED\_FIFO} de baja prioridad de tiempo real (por ejemplo,
\texttt{chrt -f 1}) y compará: ¿ahora sí lográs que los procesos
\texttt{SCHED\_OTHER} queden sin CPU? Explicá la diferencia de garantías
entre ambas políticas.
\end{extrabox}

\end{document}